\subsubsection*{2.1.7.}
对 $(\widetilde X_{(O)},Y_0)$ 的长正合列是 (2.1.5.1) 第二行的局部类似物. 对态射 $\widetilde X_{(O)}\to X_{(O)}$ 应用固有基变换, 有
\[
H^i(\widetilde X_{(O)})\xrightarrow{\ \sim\ }H^i(Y_0)
\]
故所讨论的长正合列可写为
\begin{equation}\tag{2.1.7.1}
\begin{tikzcd}[column sep=large]
H^{i-1}(X^*_{(O)}) \arrow[r]
  & H^{i-2}(Y_0)(-1) \arrow[r,"-\eta"]
  & H^i(Y_0) \arrow[r]
  & H^i(X^*_{(O)}) & {} \\
H^{i-1}(X^*) \arrow[r] \arrow[u,"\wr"]
  & H^{i-2}(Y)(-1) \arrow[r,"-\eta"] \arrow[u,equal]
  & H^i(Y) \arrow[r] \arrow[u,equal]
  & H^i(X^*) \arrow[u,dash,"\wr"]
  & {} ,
\end{tikzcd}
\end{equation}
此交换图的第二行即 (2.1.5.1) 的第二行.
\footnote{印本的标号漏去了 1 和 7 之间的点; 英译使用本节编号 (2.1.7.1).}

以下技术性引理将用于 2.2.7.

\subsubsection*{引理 2.1.8.}
\begin{center}
\begin{tikzcd}[column sep=large]
H^{n-1}(Y_0) \arrow[d,equal]
  & H^{n-1}(\widetilde X) \arrow[l,"\sim"'] \arrow[r]
  & H^{n-1}(X^*) \arrow[r]
  & H^n_{\{x\}}(X) \arrow[r]
  & H^n_c(X_s) \arrow[d,equal] \\
H^{n-1}(Y) \arrow[rrrr,"\partial"]
  & {} & {} & {} & H^n_c(X_s)
\end{tikzcd}
\end{center}
\underline{是反交换的.}\EditorialFootnote{记号核对: 原图使用 $x$ 与 $X_s$, 但本小节此前的锥记为 $X$、顶点记为 $O$, 尚未引入底 $S$. 法文原书第 181 页亦如此; 紧后的证明方块则使用 $O$ 和 $X$. 此处保留原图记号, 不将未说明的识别当作已给定.}
\footnote{印本将此引理编号为 2.7.8. 其所在位置以及 2.2.7 中的后引要求编号为 2.1.8.}

由 2.1.7, 这等价于证明下图交换:
\begin{center}
\begin{tikzcd}[column sep=large,row sep=large]
H^{n-1}(X_1-\{O\}) \arrow[r,"\partial"] \arrow[d,"\wr"]
  & H^n_{\{O\}}(X_1)=H^n_{\{O\}}(X) \arrow[d] \\
H^{n-1}(Y) \arrow[r,"\partial"]
  & H^n_c(X).
\end{tikzcd}
\end{center}
这是一般形式原理的直接结果.